\section{Formal statements and verification}\label{sec:verification}

The accompanying Lean development is authoritative for the numbered
statements and their proofs. The human-readable outlines give the
mathematical arguments and explain the hypotheses. Inputs, moves, output
and halting use the first paper's frozen machine model \cite{PaperOne}.

\paragraph{Packet coordinates.}
\lean{TapeStorageRates.PaidInputs.GrowingWidthSuccess} quantifies over a common
choice of dimension and entry width for each size threshold and accuracy
tolerance. The source machine supplies its own class of successful inputs.
\lean{TapeStorageRates.PaidClasses.compiler} realizes the eleven-tick
batch computation. \lean{TapeStorageRates.packet_code_at_threshold}
performs the finite extraction and
\lean{TapeStorageRates.family_of_packet_classes_at_threshold} chooses
dimensions to give a family. The positive lengths, two-sided degree
bounds and paired distances belong to
\lean{TapeStorageRates.SeparatedBinaryFamily}. The comparison with the
supremum of limiting upper rates is
\lean{TapeStorageRates.RateWitnesses.extremal_rates}.

\paragraph{Native coordinates.}
\lean{TapeStorageRates.NativeTypicalPartition} chooses one cut and
retained coordinate set per attained trace. Each input restricts to a
word of its own cell code. The membership proof is
\lean{TapeStorageRates.NativeTypicalPartition.projected_mem}. The tail theorem concerns the uniform
probability on $G$ under these assignments. For \cref{thm:native-tail},
depth reduction is applied to the coarse work walks; the joint support
of complete batch records pays for the physical revelations.
\Cref{cor:native} instead uses the native coding hypothesis of the
fourteenth paper at its specified physical step width and depth
\cite[Corollary~L]{PaperFourteen}.

\paragraph{Predictions.}
\lean{TapeStorageRates.PatchPredictor.missed} masks all inputs outside
a predictor's permitted helpers before evaluating its guess.
The development uses natural logarithms for Shannon entropy. Dividing
by $\ln 2$ gives the bit units of \cref{sec:patches}. Counting one scored
occurrence for each patch and target allows one target to occur in
several patches. The finite-alphabet exact-coverage result counts
distinct covered targets.

\begin{corollary}[The original qualitative target]\label{cor:formal}
\formal{TapeStorageRates/Main.lean}{TapeStorageRates.qualitative_transfer}
Under \textnormal{(T)}, for each $\varepsilon>0$ and every sufficiently large
cap there is an admissible family of limiting lower rate at least
$1-\varepsilon$.
\end{corollary}
\begin{outline}
A witness to (T) succeeds at all canonical dimensions. The original target uses $\max(1,k)$ in place of the actual tape count
$k$. The quantitative statement is proved by \lean{TapeStorageRates.quantitative_transfer}. Its qualitative consequence follows by \lean{TapeStorageRates.qualitative_of_quantitative}. At that fixed tape
count the deficit scale tends to zero. \Cref{thm:main} also gives the
bound with $k$ itself, including $k=0$.
\end{outline}

\paragraph{Reproducing the checks.}
Use the complete repository so the sibling Lean packages are available.
The project pins Lean and Mathlib v4.34.1. From this paper's directory run
\begin{verbatim}
(cd lean && lake build && lake env lean PaperAxioms.lean)
python3 tools/lean_index.py --check
\end{verbatim}
The generated axiom file queries the declarations cited in the text and
in the formal status lines. Their dependencies use at most
\texttt{propext}, \texttt{Classical.choice} and \texttt{Quot.sound}.
The finite phase calculations use locally increased heartbeat limits;
their proof terms are checked by the same kernel.
The README gives the preparation steps for a fresh checkout.
